\section{Native reconstruction inputs}
\label{app:reconstruction-inputs}
The planar and nonplanar boxes use Eqs.~\eqref{eq:planar-input} and \eqref{eq:nonplanar-input}, respectively. The high-rank targets are $(1,1,1,1,1,1,1,-u,-v)$ with $u,v\geq0$ and $1\leq u+v\leq4$. Each graph contributes fourteen training, six validation, and six reserved high-rank batches. The complete target sets differ between splits and exclude the preceding comparison batches. Individual target vectors can occur in more than one batch.

For batch index $i=0,\ldots,25$, define $j=i\bmod6$ and $q=\lfloor i/6\rfloor$. Let
\begin{align}
 (A_0,\ldots,A_5)&=(3,13,103,1009,10007,100003),\nonumber\\
 (B_0,\ldots,B_5)&=(7,19,107,1013,10009,100019).
 \label{eq:reconstruction-scales}
\end{align}
We set $s=-(A_j+2q)$ and $t=-(B_j+4q)$ in both graphs, using each graph's scalar-product convention. Indices $0$--$13$ train the model, $14$--$19$ select it, and $20$--$25$ test it. These equations define the chosen input points. In the planar graph the fourth momentum is $-p_1-p_2-p_3$, and in the nonplanar graph it is $p_1+p_2-p_3$. Expanding their squares with the listed scalar products gives zero in both cases. The test batches are specified in Table~\ref{tab:reconstruction-batches}.
\begin{table}[!htbp]\centering\small
\caption{Reserved high-rank target batches. The same $(u,v)$ sets and invariants enter both box routings. Every row defines a complete request for each graph.}
\label{tab:reconstruction-batches}
\begin{tabularx}{\linewidth}{@{}rrrY@{}}\toprule
$i$ & $s$ & $t$ & Numerator powers $(u,v)$\\\midrule
20 & -109 & -119 & $(2,0), (0,2), (0,1)$ \\
21 & -1015 & -1025 & $(2,2), (0,4)$ \\
22 & -10013 & -10021 & $(3,1), (3,0), (2,1), (1,3), (0,1)$ \\
23 & -100009 & -100031 & $(3,1), (2,2), (2,1), (2,0), (1,2), (0,4)$ \\
24 & -11 & -23 & $(2,2), (2,1), (1,2), (1,1), (0,2), (0,1)$ \\
25 & -21 & -35 & $(2,1), (1,3), (1,2), (1,0), (0,2), (0,1)$ \\
\bottomrule\end{tabularx}\end{table}

For the lower-rank box requests, the sector corner is $(1,1,1,1,1,1,1,0,0)$. A numerator request uses the supplied first-rank auxiliary targets, and a raised-power request increases one physical index. Table~\ref{tab:reconstruction-low-targets} lists their precise target vectors. Their profiles are $h=27701,27713$ for training, $27733$ for validation, and $27739$ for testing, with the low-rank kinematics of Appendix~\ref{app:families-new}.

The sunrise family has loop momenta $(\ell_1,\ell_2)$ and one external momentum $p$. Its inverse-propagator momenta and squared masses are
\begin{equation}
 (Q_1,\ldots,Q_5)=(\ell_1,\ell_2,p-\ell_1-\ell_2,p-\ell_1,p-\ell_2),
 \qquad (m_1^2,\ldots,m_5^2)=(3,7,13,0,0).
 \label{eq:reconstruction-sunrise}
\end{equation}
The first three coordinates are physical denominators and the last two are auxiliary. The three-loop vacuum uses
\begin{align}
 (Q_1,\ldots,Q_6)&=(\ell_1,\ell_2,\ell_3,\ell_1-\ell_2,\ell_2-\ell_3,\ell_1-\ell_3),\nonumber\\
 (m_1^2,\ldots,m_6^2)&=(1,1,1,13,13,2h+53).
 \label{eq:reconstruction-vacuum}
\end{align}
All six coordinates are physical. Both structures use the same profile split as the low-rank boxes. For the sunrise, $p^2=-1009-12h$. These assignments specify test families, rather than universal relations between masses and invariants. In every native request the kinematic quantities are fixed rational numbers and $d$ remains symbolic, so the coefficients belong to $\Q(d)$. This test does not reconstruct general symbolic scattering invariants.

\begin{table}[!htbp]\centering\small
\caption{Lower-rank reserved targets. Training and validation use the same index patterns at different physical points. A corner request contains the all-unit physical sector.}
\label{tab:reconstruction-low-targets}
\begin{tabularx}{\linewidth}{@{}lY@{}}\toprule
Structure and request & Target indices\\\midrule
Planar box, corner & $(1,1,1,1,1,1,1,0,0)$ \\
Planar box, numerator & $(1,1,1,1,1,1,1,0,-1), (1,1,1,1,1,1,1,-1,0)$ \\
Planar box, raised power & $(1,1,1,1,1,2,1,0,0), (1,1,1,1,1,1,2,0,0), (2,1,1,1,1,1,1,0,0)$ \\
Nonplanar box, corner & $(1,1,1,1,1,1,1,0,0)$ \\
Nonplanar box, numerator & $(1,1,1,1,1,1,1,0,-1), (1,1,1,1,1,1,1,-1,0)$ \\
Nonplanar box, raised power & $(1,1,1,1,1,2,1,0,0), (1,1,1,1,1,1,2,0,0), (2,1,1,1,1,1,1,0,0)$ \\
Sunrise, corner & $(1,1,1,0,0)$ \\
Sunrise, raised power & $(1,2,1,0,0), (1,1,2,0,0), (2,1,1,0,0)$ \\
Vacuum, corner & $(1,1,1,1,1,1)$ \\
Vacuum, raised power & $(1,2,1,1,1,1), (1,1,2,1,1,1), (1,1,1,2,1,1)$ \\
\bottomrule\end{tabularx}\end{table}

The protected native bases consist of the integral indices below. Each listed vector $\boldsymbol b$ denotes $M_{\boldsymbol b}=I_{\boldsymbol b}$ in its corresponding family. The independently prepared symbolic and reconstructed calculations use the same basis for each request. The lists specify the comparison basis; they do not assert a proof of global master independence from a finite seed system.

Planar box:\[\begin{aligned}\{&(0,0,1,1,0,0,1,0,0),(0,1,0,0,0,1,1,0,0),\\&(0,1,0,1,1,0,1,0,0),(0,1,1,1,0,1,1,0,0),\\&(1,0,1,1,1,0,0,0,0),(1,1,1,0,0,1,1,0,0),\\&(1,1,1,1,1,1,1,-1,0),(1,1,1,1,1,1,1,0,0)\}.\end{aligned}\]

Nonplanar box:\[\begin{aligned}\{&(0,0,1,1,1,1,0,0,0),(0,1,0,0,1,1,0,0,0),\\&(0,1,1,1,1,1,1,0,0),(1,0,0,1,0,1,0,0,0),\\&(1,0,1,1,1,1,0,0,0),(1,1,0,0,0,0,1,0,0),\\&(1,1,0,1,0,1,1,0,0),(1,1,0,1,1,1,0,0,0),\\&(1,1,1,0,1,0,1,0,0),(1,1,1,1,0,0,1,0,0),\\&(1,1,1,1,1,1,1,-1,0),(1,1,1,1,1,1,1,0,0)\}.\end{aligned}\]

Sunrise, basis 1:\[\begin{aligned}\{&(0,1,1,0,0),(1,0,1,0,0),\\&(1,1,0,0,0),(1,1,1,-1,0),\\&(1,1,1,0,-1),(1,1,1,0,0),\\&(1,1,2,0,0)\}.\end{aligned}\]

Sunrise, basis 2:\[\begin{aligned}\{&(1,1,1,0,0)\}.\end{aligned}\]

Three-loop vacuum, basis 1:\[\begin{aligned}\{&(0,1,1,1,0,1),(0,1,1,1,0,2),\\&(0,1,1,2,0,1),(1,0,0,1,1,1),\\&(1,0,1,1,0,1),(1,0,1,1,1,0),\\&(1,0,1,1,1,1),(1,0,1,1,2,0),\\&(1,1,0,1,0,1),(1,1,0,1,1,0),\\&(1,1,0,1,1,1),(1,1,1,0,0,1),\\&(1,1,1,1,0,0),(1,1,1,1,0,1),\\&(1,1,1,1,1,0),(1,1,1,1,1,1)\}.\end{aligned}\]

Three-loop vacuum, basis 2:\[\begin{aligned}\{&(1,1,1,1,1,1)\}.\end{aligned}\]


\section{Coefficient reconstruction}
\label{app:reconstruction-algebra}
Kira and FireFly provide the finite-field algorithms used here~\cite{Kira2,FireFly2019,FireFly2020}. To explain the relation to an exact rational coefficient, suppose
\begin{equation}
 r(d)=P(d)/Q(d),\qquad
 P(d)=\sum_{j=0}^{n}p_jd^j,\quad Q(d)=\sum_{j=0}^{m}q_jd^j.
 \label{eq:reconstruction-rational}
\end{equation}
One nonzero denominator coefficient fixes the overall normalization. At a nonsingular point $d=z$ modulo a prime, a numerical reduction returns $r(z)$. Multiplying by the denominator gives the linear constraint
\begin{equation}
 \sum_{j=0}^{n}p_jz^j-r(z)\sum_{j=0}^{m}q_jz^j=0\pmod p.
 \label{eq:reconstruction-sample}
\end{equation}
This follows directly from Eq.~\eqref{eq:reconstruction-rational}. With correct degree bounds and sufficiently many independent nonsingular samples, the constraints determine the normalized coefficients over the field. Combining suitable primes and reconstructing rational numbers then supplies rational coefficients~\cite{FireFly2019,FireFly2020}. FireFly determines the necessary information adaptively; our model chooses the native procedure and probe batch size.

The field operations are exact, while reconstruction termination relies on its degree discovery and checks. Our evaluation therefore also subtracts every reconstructed coefficient from a separately computed symbolic reference and requires the rational difference to vanish. This checks the actual tested answer. Neither a model score nor a successful reduction to the supplied basis alone proves every coefficient for an arbitrary unseen request.
